% Introduction -- written from scratch on 2026-10-02 after the 2026-09-07 rejection.
% Rewritten 2026-10-04 after the funnel of thesis/style_guide_aida.md (her ch. 1): one step per
% paragraph, claims only. Paragraph 1 is fixed. Evidence cut from here and still owed to the
% State of the Art: 196 patients, left main bifurcation angle 35.5 to 178 degrees
% \cite{temov2016bifurcation}; 136 autopsy hearts, no visible narrowing until plaque fills about 40\%
% of the cross-section \cite{glagov1987remodeling}. The 127-patient curvature result is already in
% clinical_background/05_hemodynamics.tex.
% 2026-10-04: limitations paragraph trimmed to a problem statement; the evidence (Han, Tello Ayala,
% Glagov, Temov) now lives in literature_review.tex. "considers only" softened to "rests on", since
% chan2024orfan shows imaging already refines risk in referred patients. shen2026geometry replaced
% by kashyap2022tortuosity and zebicmihic2024index (Shen excluded thesis-wide, user 2026-10-04).
% 2026-10-05: paragraph 4 reordered to match the review (measurement = Sec. 2, timing = Sec. 3);
% zebicmihic2024index swapped for cui2017bifurcation so every cited study is discussed in the review.
% No forward references in this chapter (user).
% 2026-11-14 (user): paragraph 3 names SCORE2 and that it gives the 10-year risk of fatal and non-fatal
%   CVD (score2_2021, Methods and Supplementary Box 1), so the reader knows it is a 10-year probability
%   before sec:tech-recalibration.
% 2026-11-15 (user): paragraphs 5 and 6 rewritten to the current design and the prognostic factor
%   framing (root CLAUDE.md): association beyond the SCORE2 risk factors, plaque-free baseline,
%   in-house tracer, four features, one pooled segment-level model, one test, prediction secondary.
%   Four features, not the five of the root CLAUDE.md: the Methods notes (04_model_size.tex,
%   05c_daughter_ratio.tex, 2026-11-13) replaced the LM->LAD and LM->LCx ratios by the daughter
%   asymmetry ln(D_LCx/D_LAD) and dropped Finet's ratio.
%   Cut: own segmentation pipeline, logistic regression + transformer, per-artery models, "new
%   plaque", "identify individuals at increased risk" (a risk-model claim).
%   The words "prognostic factor" are deliberately not used: the term is open (ask supervisor).
%   The model is not named (cloglog GLMM, LR test); that is left to the Technical Background and
%   the Methods. CHECK once progression_who_pooled.py is in thesis_final/: the model and test
%   sentences still match the code.
%   CHECK: "four features" and "narrowing of their lumen" (taper per artery or per segment is
%   unresolved in 05d_taper.tex).
% 2026-11-15 (user): the association is the objective from the start, not a change of plan.
%   Paragraph 4: the gap is stated as the association beyond the risk factors, not "adds to such a
%   score" (a risk-model question); "a consequence of the disease rather than its cause" replaced by
%   "the disease may have altered the geometry that was measured" (no causal wording).
%   Paragraph 6: the opener "not as a risk model for clinical use and not as a cause of plaque" cut;
%   the causal limit is now its own sentence with its reason (adjustment only for the risk factors).
% 2026-11-15 (user): base decided (root CLAUDE.md): the six SCORE2 risk factors as separate
%   covariates, SCORE2 only as the reason for the set. The chapter is no longer centred on the score.
%   Paragraph 3: "combining ... into a single score such as SCORE2, which gives the 10-year risk of
%   fatal and non-fatal cardiovascular disease" replaced by "Risk scores such as SCORE2 are built
%   from age, sex, smoking, blood pressure and cholesterol"; SCORE2 is now named once, as the source
%   of the factors. Cut: the 10-year meaning of the score (added 2026-11-14). It is owed to where the
%   recalibrated score is described as a sensitivity analysis (appendix/score2_recalibration.tex).
%   Paragraph 5: "and to the risk factors of SCORE2" -> "and to the established risk factors", as in
%   the aim sentence; the TODO on the open base removed. The sentence does not say how the factors
%   enter the model, so it holds while progression_who_pooled.py still defaults to log SCORE2.
%   CHECK: once that code is in thesis_final/ with the six factors as its default.

\chapter{Introduction}
\label{ch:introduction}
Cardiovascular disease is the leading cause of death worldwide, accounting for an estimated 19.2
million deaths in 2023 \cite{stark2025gbd}.
Coronary artery disease (CAD) is the largest single contributor, responsible for nearly half of
these deaths \cite{stark2025gbd}.
It develops when plaque builds up in the walls of the coronary arteries, which supply the heart
muscle with blood.
The plaque slowly narrows the arteries and can, in extreme cases, rupture, forming a clot that
blocks blood flow and causes a heart attack \cite{libby2019atherosclerosis}.

Plaque does not form evenly along the coronary tree, even though risk factors such as high blood
pressure, high cholesterol and smoking reach every segment through the blood.
Instead, it tends to collect where the drag of the blood flow on the vessel wall is low or
oscillating \cite{chatzizisis2007ess}.
This drag is known as wall shear stress, and it is low mainly on the outer wall just after a
bifurcation and on the inner wall of a bend \cite{asakura1990flow}.
Where plaque forms is therefore tied to the geometry of the tree, which, unlike shear, can be
measured directly from an image.

Coronary geometry varies considerably between individuals \cite{temov2016bifurcation}, therefore some
trees contain more extensive low-shear regions than others.
Predisposition to plaque may therefore depend not only on systemic risk factors but also on the
anatomy of the individual's own arteries.
Standard risk assessment, however, rests only on the former.
Risk scores such as SCORE2 are built from age, sex, smoking, blood pressure and cholesterol
\cite{score2_2021}.

Whether geometry is associated with developing plaque beyond these risk factors remains unknown.
Studies relating geometry to plaque have defined geometry inconsistently, and the definition can
decide whether a relationship is found \cite{cui2017bifurcation,kashyap2022tortuosity}.
Most also measured geometry in referred patients whose plaque was already present, so the disease
may have altered the geometry that was quantified \cite{han2022plaque}.

The aim of this thesis is to test whether coronary geometry, measured before any plaque is visible,
is associated with developing plaque beyond the established risk factors.
Rather than referred patients, this work uses adults selected at random from the population of
greater Copenhagen by the Copenhagen General Population Study (CGPS) and scanned for research with
coronary computed tomography angiography (CCTA) \cite{fuchs2023cgps}.
A longitudinal subset of these individuals was scanned a second time about 10 years later.
In these scans, the centerline and the lumen of the three main coronary arteries are extracted with
an in-house tracer, and four features are computed from them: the curvature of the arteries, the
narrowing of their lumen along their length, the angle at which the left main artery divides, and
the difference in diameter between its two daughter arteries.
A single regression model then relates plaque development in each artery segment to these features
and to the risk factors, with all segments of an individual analysed together.
The association is tested once, by comparing this model with the same model without the geometry.
How much geometry adds to the prediction of who develops plaque is estimated as a secondary result.

An association would indicate that the anatomy of an individual's arteries carries information
about who develops plaque that the established risk factors do not.
It would not show that geometry causes plaque, since the analysis adjusts only for these
factors.
The absence of an association would suggest that geometry, as it can be measured from CCTA, adds
little to them, even though the location of plaque within a tree follows the geometry.
